\section{Anhang}
\subsection{Approfondimento taylor Reihe}

Questo approccio permette di sviluppare serie di Taylor per funzioni razionali evitando
il laborioso calcolo di  utilizzando la serie 
geometrica.

\begin{itemize}
    \item \textbf{1. Partialbruchzerlegung (PBZ):} Si scompone la funzione frazionaria complessa
    \begin{equation*}
        f(x) = \frac{P(x)}{Q(x)} \xrightarrow{PBZ} \frac{A}{ax+b} + \frac{B}{cx+d}
    \end{equation*}

    \item \textbf{2. Geometrische Reihe:} 
    si riconducono i termini alla forma chiusa della serie geometrica:
    \begin{equation*}
        \frac{1}{1-\cor{q}} = \sum_{n=0}^{\infty} \cor{q}^n
    \end{equation*}

    \item \textbf{3. Normalizzazione del Denominatore:} Si manipolano i fratti 
    raccogliendo le costanti a denominatore per ottenere esattamente un $1$ iniziale, forma standard $\frac{1}{1-q}$.
    \begin{equation*}
        \frac{A}{a+bx} = \frac{A}{a} \cdot \frac{1}{1 - \left(-\frac{b}{a}x\right)} \quad \implies \quad \cor{q} = -\frac{b}{a}x
    \end{equation*}

    \item \textbf{4. Linearità e Vorfaktor:} Le costanti e segni vengono spostati nelle 
    sommatorie. Si uniscono le serie in un'unica sommatoria raccogliendo $x^n$.

    \item \textbf{5. Konvergenzradius ($\rho$):} La serie geometrica converge solo se 
    $|q| < 1$. 
    
    Se la $f$ è composta da più serie, si calcola la condizione di convergenza 
    per ognuna. Il raggio di convergenza totale è l'intersezione delle condizioni 
    (\textbf{minimo} raggi)
\end{itemize}

\subsection{Cambio di indice nelle sommatorie}

Per traslare l'indice di una sommatoria di una quantità costante $s$, si applica la seguente trasformazione generale:

\begin{equation*}
\sum_{n=n_0}^{N} a_n \quad \xrightarrow{k = n + s} \quad \sum_{k=n_0+s}^{N+s} a_{k-s}
\end{equation*}

I passaggi logici sono tre:
\begin{enumerate}
    \item \textbf{Nuova variabile}: Si definisce $k = n + s$.
    \item \textbf{Nuovi limiti}: inferiore diventa $k_{start} = n_0 + s$ e superiore $k_{end} = N + s$.
    \item \textbf{Sostituzione}: $n$ in funzione di $k$ ($n = k - s$) e sostituire all'interno dell'argomento.
\end{enumerate}

\begin{equation*}
\sum_{n=0}^{\infty} \frac{x^{n+1}}{(n+1)^2} \quad \xrightarrow{k = n + 1} \quad \sum_{k=1}^{\infty} \frac{x^k}{k^2}
\end{equation*}
In questo caso, avendo $n = k - 1$, il termine $(n+1)$ diventa semplicemente $k$.

\subsection{Ordine di contatto $n$}

Per determinare l'ordine di contatto $n$ in un punto $x_0$ tra due funzioni $f(x)$ e $g(x)$, è possibile sviluppare entrambe in \textbf{serie di Taylor} (o serie di potenze) e confrontare i rispettivi coefficienti $a_k$.

Due funzioni hanno un contatto di ordine $n$ se:

i primi $n+1$ coefficienti delle loro serie sono identici.

\textbf{Nota bene:} 
\begin{itemize}
    \item I coefficienti nulli (\(a_k = 0\)) sono fondamentali e vanno conteggiati nel confronto.
    \item Il primo coefficiente differente determina dove il contatto si "rompe".
\end{itemize}

\subsection{Singularitätstest ergänzung}\label{chap:sing_erg}
si dividono spesso i membri dell'equazione per termini contenenti la variabile dipendente. 
Questo processo esclude implicitamente le soluzioni che annullano tali termini.

\subsubsection{Procedimento Generale}

\begin{enumerate}
    \item \textbf{Identificazione dei punti critici:} 
    Individua ogni termine $g(y)$ o $f(z)$ che è stato posto a denominatore durante i passaggi algebrici.
    
    \item \textbf{Ricerca dei candidati:}
    Imponi il termine uguale a zero per trovare potenziali soluzioni 
    costanti o funzioni particolari: $g(y) = 0 \implies y = c $
    
    \item \textbf{DGL (Einsetzen):}
    Sostituisci la funzione candidata $y$ e la sua derivata $y'$ nella DGL iniziale
    
    \item \textbf{Valutazione del risultato:}
    \begin{itemize}
        \item \textbf{Soluzione:} identità sempre vera (es. $0 = 0$), \cgn{la funzione è una soluzione valida.} 
        \item \textbf{Test Fallito:} un'uguaglianza dipendente da $x$ (es. $-x^2 = 0$) la funzione \underline{non} è una soluzione della DGL.
    \end{itemize}
\end{enumerate}

\subsection{Neutralità del Faltungsintegral}
Se si calcola la soluzione particolare $y_p(x)$ tramite il \textbf{Faltungsintegral} con estremo inferiore uguale a $0$.
per definizione geometrica e analitica si ha sempre:
\[ y_p(0) = 0 \quad \text{e} \quad y_p'(0) = 0 \]

Se le anfangsbedingungen hanno $\mathbf{x_0 = 0}$, $y_p$ è totalmente \textbf{neutra} nel calcolo delle costanti $A$ e $B$.

\cor{Si possono quindi ricavare $A$ e $B$ imponendo le condizioni \textbf{solo} sulla parte omogenea}